\documentclass[11pt]{article}

\usepackage[T1]{fontenc}
\usepackage[utf8]{inputenc}
\usepackage[margin=1in]{geometry}
\usepackage{amsmath,amssymb}
\usepackage{xcolor}
\usepackage{hyperref}
\usepackage{listings}

\hypersetup{
    colorlinks=true,
    linkcolor=blue!50!black,
    urlcolor=blue!50!black
}

\lstdefinestyle{codeblock}{
    basicstyle=\ttfamily\small,
    breaklines=true,
    columns=fullflexible,
    keepspaces=true,
    showstringspaces=false,
    frame=single,
    framerule=0.3pt,
    rulecolor=\color{black!30},
    backgroundcolor=\color{black!2}
}

\newcommand{\code}[1]{\texttt{\detokenize{#1}}}
\newcommand{\dtau}{\Delta\tau}
\newcommand{\ET}{E_{\mathrm{T}}}
\newcommand{\EL}{E_{\mathrm{L}}}

\title{GFMC in This Repository: A Theory-to-Code Walkthrough}
\author{}
\date{\today}

\begin{document}

\maketitle
\tableofcontents

\section{Purpose}

This note explains how the repository's GFMC implementation realizes the
imaginary-time projection idea
\begin{equation}
\Psi_{n+1} = e^{-\dtau(H-\ET)} \Psi_n,
\end{equation}
with direct reference to the actual Julia code under
\code{src/UseCases/gfmc/}.

The main point is that this code is \emph{not} a branching-population DMC
engine. It is a \emph{fixed-population, weighted GFMC projector}:
\begin{itemize}
    \item walker coordinates are stored in a dense matrix,
    \item each walker carries a multiplicative weight,
    \item those weights are periodically resampled back to an equally weighted
    ensemble,
    \item optional importance sampling and optional fixed-node killing modify
    the short-time kernel.
\end{itemize}

So the clean conceptual summary is:
\begin{quote}
This code applies a short-time approximation to the imaginary-time Green's
function, samples the move part with Gaussian or drift-diffusion proposals,
accumulates the diagonal part into walker weights, and periodically
reconfigures to keep exactly \code{targetN} walkers.
\end{quote}

\section{What Object the Code Represents}

At the operator level, GFMC repeatedly applies
\begin{equation}
G = e^{-\dtau(H-\ET)}.
\end{equation}
In coordinate space this is the integral transform
\begin{equation}
\Psi_{n+1}(R) = \int dR' \, G(R,R';\dtau)\Psi_n(R').
\end{equation}

With importance sampling by a trial state $\Psi_T$, it is more natural to evolve
\begin{equation}
f(R,\tau) = \Psi_T(R)\Psi(R,\tau),
\end{equation}
using an importance-sampled kernel $\widetilde G$.

The code stores a weighted ensemble
\begin{equation}
f_n(R) \approx \sum_{j=1}^{N_w} w_j \, \delta(R-R_j),
\end{equation}
where:
\begin{itemize}
    \item the columns $R_j$ live in \code{sim.positions},
    \item the weights $w_j$ live in \code{sim.weights},
    \item the batch data \code{drift}, \code{logpsi}, \code{branch\_energy},
    and \code{sign} hold the local information needed to apply the short-time
    kernel.
\end{itemize}

That is the key bridge from the math to the implementation: the integral
operator is represented by repeated stochastic updates of
\code{(positions, weights)}.

\section{Core State and Interfaces}

\subsection{Parameters and per-walker data}

The runtime parameters are deliberately small:

\begin{lstlisting}[style=codeblock,caption={\code{src/UseCases/gfmc/params.jl}}]
struct GFMCParams
    dt::Float64
    nsteps::Int
    nequil::Int
    targetN::Int
    ET0::Float64
    feedback::Float64
    reconfiguration_interval::Int
    branch_cap::Float64
    energy_window::Int
end
\end{lstlisting}

These map almost one-to-one onto the projector formula:
\begin{itemize}
    \item \code{dt} is $\dtau$,
    \item \code{ET0} is the initial reference energy $\ET$,
    \item \code{reconfiguration\_interval} determines how often the weighted
    ensemble is resampled,
    \item \code{branch\_cap} limits very large short-time weights.
\end{itemize}

The auxiliary data carried for every walker is:

\begin{lstlisting}[style=codeblock,caption={\code{src/UseCases/gfmc/kernels.jl}}]
mutable struct GFMCBatchData{M<:AbstractMatrix{Float64},V<:AbstractVector{Float64}}
    drift::M
    logpsi::V
    branch_energy::V
    sign::V
end
\end{lstlisting}

This already tells you the structure of the implemented kernel:
\begin{itemize}
    \item \code{drift} defines the proposal mean for guided moves,
    \item \code{logpsi} enters the Metropolis correction,
    \item \code{branch\_energy} provides the exponential weight factor,
    \item \code{sign} is the fixed-node diagnostic.
\end{itemize}

\subsection{The simulation container}

The full GFMC state is stored in \code{GFMCSim}:

\begin{lstlisting}[style=codeblock,caption={Excerpt from \code{src/UseCases/gfmc/sim.jl}}]
mutable struct GFMCSim{
    K<:AbstractGFMCKernel,
    RNG,
    B<:AbstractGFMCBackend,
    R<:AbstractGFMCReconfiguration,
    M<:AbstractMatrix{Float64},
    V<:AbstractVector{Float64},
}
    kernel::K
    params::GFMCParams
    positions::M
    proposal_positions::M
    state_data::GFMCBatchData{M,V}
    proposal_data::GFMCBatchData{M,V}
    weights::V
    rng::RNG
    backend::B
    reconfiguration::R
    step::Int
    ET::Float64
    ET_history::Vector{Float64}
    population_history::Vector{Int}
    energy_mean_history::Vector{Float64}
    energy_variance_history::Vector{Float64}
    effective_population_history::Vector{Float64}
    mean_weight_history::Vector{Float64}
    acceptance_history::Vector{Float64}
    walker_positions_history::Vector{Vector{Vector{Float64}}}
end
\end{lstlisting}

There are two important design choices here:
\begin{enumerate}
    \item Walkers are stored column-wise in a matrix, so one GFMC step can act
    on a whole batch.
    \item The code keeps both current and proposal buffers, then swaps them
    after a step instead of reallocating.
\end{enumerate}

\section{Kernel Abstraction}

The GFMC engine itself does not know the physics. It asks the kernel object for
the information needed to build a short-time propagator:

\begin{lstlisting}[style=codeblock,caption={Excerpt from \code{src/UseCases/gfmc/abstract.jl}}]
abstract type AbstractGFMCKernel end
abstract type AbstractGFMCBackend end
abstract type AbstractGFMCReconfiguration end

configuration_dimension(::AbstractGFMCKernel) = error("configuration_dimension not implemented")
diffusion_constant(::AbstractGFMCKernel) = error("diffusion_constant not implemented")
boundary(::AbstractGFMCKernel) = error("boundary not implemented")
uses_importance_sampling(::AbstractGFMCKernel) = false
uses_fixed_node(::AbstractGFMCKernel) = false
\end{lstlisting}

This is the code form of the abstract GFMC statement
\begin{equation}
f_{n+1} = \widetilde G f_n.
\end{equation}
The engine knows \emph{how} to apply a short-time kernel, while the kernel type
supplies \emph{what} drift, local energy, sign, diffusion constant, and
boundary geometry mean for a particular Hamiltonian.

\subsection{Generic Hamiltonian path}

For a generic Hamiltonian, the importance-sampled quantities are filled by the
routine shown below:

\begin{lstlisting}[style=codeblock,caption={Generic kernel evaluation in \code{src/UseCases/gfmc/kernels.jl}}]
function evaluate_configuration_data!(kernel::GenericGFMCKernel, data::GFMCBatchData, X::AbstractMatrix{<:Real})
    uses_guiding = uses_importance_sampling(kernel)
    @inbounds for j in axes(X, 2)
        R = @view X[:, j]
        if uses_guiding
            drift_j = drift(kernel.guiding, R)
            for i in axes(X, 1)
                data.drift[i, j] = Float64(drift_j[i])
            end
            data.logpsi[j] = Float64(kernel.guiding.trial.logpsi(R))
            data.branch_energy[j] = Float64(local_energy(kernel.guiding, R))
            data.sign[j] = Float64(signpsi(kernel.guiding, R))
        else
            for i in axes(X, 1)
                data.drift[i, j] = 0.0
            end
            data.logpsi[j] = 0.0
            data.branch_energy[j] = Float64(potential(kernel.H, R))
            data.sign[j] = 1.0
        end
    end
    return data
end
\end{lstlisting}

So in the guided case, the code is exactly using the standard DMC/GFMC objects
\begin{align}
F(R) &= 2\nabla \log |\Psi_T(R)|, \\
\EL(R) &= -D\left(\nabla^2 \log |\Psi_T(R)| + |\nabla \log |\Psi_T(R)||^2\right) + V(R).
\end{align}

This comes from the shared guiding helper:

\begin{lstlisting}[style=codeblock,caption={\code{src/Domain/guiding.jl}}]
drift(g::ImportanceGuiding, R) = 2 .* g.trial.gradlogpsi(R)

function local_energy(g::ImportanceGuiding, R)
    D = diffusion_constant(g.H)
    gradlog = g.trial.gradlogpsi(R)
    lapllog = g.trial.lapllogpsi(R)
    V = potential(g.H, R)
    return -D * (lapllog + sum(abs2, gradlog)) + V
end
\end{lstlisting}

\subsection{Specialized fermion-ring path}

The repository also includes a hand-written kernel for
\code{SpinlessFermionRing1D}. The crucial excerpt is:

\begin{lstlisting}[style=codeblock,caption={Excerpt from \code{src/UseCases/gfmc/kernels.jl}}]
if kernel.use_guiding
    for i in 1:model.N
        xi = X[i, j]
        coskx = cos(model.k_lat * xi)
        sinkx = sin(model.k_lat * xi)
        logpsi_j += model.lambda * coskx
        data.drift[i, j] = -2.0 * model.lambda * model.k_lat * sinkx
        lapl_j += -model.lambda * model.k_lat^2 * coskx
        potential_j += model.V0 * coskx
    end

    for i in 1:(model.N - 1)
        for k in (i + 1):model.N
            dx_ik = displacement(model.bc, X[k, j], X[i, j])
            u = model.alpha_pair * dx_ik
            s = sin(u)
            abs_s = abs(s)

            if abs_s < model.node_tol
                sign_j = 0.0
            elseif sign_j != 0.0 && s < 0
                sign_j = -sign_j
            end

            s_eff = abs_s < model.trig_eps ? (s >= 0 ? model.trig_eps : -model.trig_eps) : s
            logpsi_j += log(abs(s_eff))

            c = model.alpha_pair * cos(u) / s_eff
            data.drift[i, j] += 2.0 * c
            data.drift[k, j] -= 2.0 * c

            lapl_j += -2.0 * (model.alpha_pair^2) / (s_eff * s_eff)
        end
    end

    grad_norm2 = 0.0
    for i in 1:model.N
        grad_i = 0.5 * data.drift[i, j]
        grad_norm2 += grad_i * grad_i
    end

    data.logpsi[j] = logpsi_j
    data.branch_energy[j] = -model.D * (lapl_j + grad_norm2) + potential_j
    data.sign[j] = sign_j
end
\end{lstlisting}

This is where the model-specific trial state is encoded:
\begin{itemize}
    \item the one-body lattice factor contributes the $\lambda \cos(k_{\rm lat}x_i)$
    term,
    \item the pairwise $\sin(\alpha_{\rm pair} \Delta x_{ik})$ factor carries
    antisymmetry,
    \item \code{sign\_j} is computed from that pair structure, so fixed-node
    killing has real content here.
\end{itemize}

\section{The GFMC Step as Implemented}

At the mathematical level, the code is using a short-time kernel of the form
\begin{equation}
\widetilde G(R',R;\dtau)
\approx
T(R' \leftarrow R)
\exp\left[
-\frac{\dtau}{2}\bigl(\EL(R')+\EL(R)-2\ET\bigr)
\right],
\end{equation}
where $T$ is either a plain Gaussian diffusion proposal or a drifted Gaussian
proposal corrected by Metropolis-Hastings.

\subsection{Proposal generation}

The first part of \code{step!} is:

\begin{lstlisting}[style=codeblock,caption={Proposal generation in \code{src/UseCases/gfmc/run.jl}}]
function step!(sim::GFMCSim)
    D = diffusion_constant(sim.kernel)
    dt = sim.params.dt
    sigma = sqrt(2.0 * D * dt)
    nwalkers = size(sim.positions, 2)
    bc = boundary(sim.kernel)

    if uses_importance_sampling(sim.kernel)
        randn!(sim.backend, sim.rng, sim.proposal_positions)
        @inbounds for j in 1:nwalkers
            for i in axes(sim.positions, 1)
                sim.proposal_positions[i, j] =
                    sim.positions[i, j] +
                    D * dt * sim.state_data.drift[i, j] +
                    sigma * sim.proposal_positions[i, j]
            end
        end
    else
        randn!(sim.backend, sim.rng, sim.proposal_positions)
        @inbounds for j in 1:nwalkers
            for i in axes(sim.positions, 1)
                sim.proposal_positions[i, j] =
                    sim.positions[i, j] + sigma * sim.proposal_positions[i, j]
            end
        end
    end

    wrap_configurations!(sim.kernel, sim.proposal_positions)
    evaluate_configuration_data!(sim.kernel, sim.proposal_data, sim.proposal_positions)
\end{lstlisting}

This is the code version of:
\begin{align}
R' &= R + \sqrt{2D\dtau}\,\eta && \text{(unguided)}, \\
R' &= R + D\dtau\,F(R) + \sqrt{2D\dtau}\,\eta && \text{(guided)}.
\end{align}

The \code{wrap\_configurations!} call is important: because the engine supports
periodic boundaries, all proposals are mapped back into the simulation cell
before their local data is evaluated.

\subsection{Fixed-node killing, acceptance, and branching}

The heart of the step is the walker loop:

\begin{lstlisting}[style=codeblock,caption={Walker update logic in \code{src/UseCases/gfmc/run.jl}}]
@inbounds for j in 1:nwalkers
    killed = false

    if uses_fixed_node(sim.kernel)
        sold = sim.state_data.sign[j]
        snew = sim.proposal_data.sign[j]
        killed = (sold == 0.0) || (snew == 0.0) || (sold != snew)
    end

    if killed
        for i in axes(sim.positions, 1)
            sim.proposal_positions[i, j] = sim.positions[i, j]
            sim.proposal_data.drift[i, j] = sim.state_data.drift[i, j]
        end
        sim.proposal_data.logpsi[j] = sim.state_data.logpsi[j]
        sim.proposal_data.branch_energy[j] = sim.state_data.branch_energy[j]
        sim.proposal_data.sign[j] = sim.state_data.sign[j]
        sim.weights[j] = 0.0
        continue
    end

    if uses_importance_sampling(sim.kernel)
        denom = 4.0 * D * dt
        sq_f = 0.0
        sq_b = 0.0

        for i in axes(sim.positions, 1)
            mu_f = sim.positions[i, j] + D * dt * sim.state_data.drift[i, j]
            mu_b = sim.proposal_positions[i, j] + D * dt * sim.proposal_data.drift[i, j]
            df = displacement(bc, mu_f, sim.proposal_positions[i, j])
            db = displacement(bc, mu_b, sim.positions[i, j])
            sq_f += df * df
            sq_b += db * db
        end

        log_ratio = -(sq_b - sq_f) / denom +
                    2.0 * (sim.proposal_data.logpsi[j] - sim.state_data.logpsi[j])
        if log(rand(sim.rng)) >= min(0.0, log_ratio)
            for i in axes(sim.positions, 1)
                sim.proposal_positions[i, j] = sim.positions[i, j]
                sim.proposal_data.drift[i, j] = sim.state_data.drift[i, j]
            end
            sim.proposal_data.logpsi[j] = sim.state_data.logpsi[j]
            sim.proposal_data.branch_energy[j] = sim.state_data.branch_energy[j]
            sim.proposal_data.sign[j] = sim.state_data.sign[j]
        else
            accept_count += 1
        end
    else
        accept_count += 1
    end

    sim.weights[j] *= _branch_weight(
        sim,
        sim.state_data.branch_energy[j],
        sim.proposal_data.branch_energy[j],
    )
end
\end{lstlisting}

There are three separate GFMC ideas embedded here.

\paragraph{1. Fixed-node constraint.}
If the old and new signs disagree, or if either sign is zero, the walker is not
allowed to cross the node. In this implementation the walker is reverted and
its weight is set to zero. So fixed-node here is realized as \emph{killing plus
later resampling}, not by reflecting the move.

\paragraph{2. Importance-sampled acceptance.}
The forward and backward drifted Gaussians are compared using the
Metropolis-Hastings ratio
\begin{equation}
\log A =
-\frac{1}{4D\dtau}\left(\|R-\mu_b\|^2-\|R'-\mu_f\|^2\right)
+ 2\bigl(\log|\Psi_T(R')|-\log|\Psi_T(R)|\bigr),
\end{equation}
with periodic-boundary-aware displacements from
\code{src/Domain/boundaries1D.jl}. This is exactly what enforces detailed
balance for the drifted proposal.

\paragraph{3. Short-time branching weight.}
The multiplicative weight factor is defined by:

\begin{lstlisting}[style=codeblock,caption={\code{src/UseCases/gfmc/run.jl}}]
function _branch_weight(sim::GFMCSim, Eold::Real, Enew::Real)
    logw = -0.5 * sim.params.dt * ((Float64(Eold) + Float64(Enew)) - 2.0 * sim.ET)
    w = exp(clamp(logw, -700.0, 700.0))
    return min(w, sim.params.branch_cap)
end
\end{lstlisting}

This is the trapezoidal short-time approximation
\begin{equation}
w_{\rm branch} \approx
\exp\left[
-\frac{\dtau}{2}\bigl(\EL(R_{\rm old})+\EL(R_{\rm new})-2\ET\bigr)
\right],
\end{equation}
with two practical safeguards:
\begin{itemize}
    \item the logarithm is clamped before exponentiation,
    \item the resulting weight is capped by \code{branch\_cap}.
\end{itemize}

\subsection{Buffer swap and diagnostics}

After all walkers are processed, the current and proposal buffers are swapped.
Then the code records
\begin{equation}
N_{\rm eff} = \frac{\left(\sum_j w_j\right)^2}{\sum_j w_j^2},
\end{equation}
the mean weight, and the acceptance rate. This is the right diagnostic because
the actual number of stored walkers remains fixed even when the weighted
ensemble has effectively collapsed.

\section{Reconfiguration: Fixed Population Instead of Branching}

The defining difference from branching DMC appears in
\code{src/UseCases/gfmc/reconfiguration.jl} and \code{src/UseCases/gfmc/run.jl}.

First, weights are normalized into probabilities:

\begin{lstlisting}[style=codeblock,caption={\code{src/UseCases/gfmc/reconfiguration.jl}}]
function _normalized_weights(weights::AbstractVector{<:Real})
    total = sum(weights)
    isfinite(total) && total > 0 || error("All GFMC walker weights vanished.")
    probs = Vector{Float64}(undef, length(weights))
    @inbounds for i in eachindex(weights)
        wi = Float64(weights[i])
        wi >= 0 || throw(ArgumentError("GFMC weights must be non-negative"))
        probs[i] = wi / total
    end
    return probs, total / length(weights)
end
\end{lstlisting}

Then reconfiguration chooses indices and resets the ensemble to equal weights:

\begin{lstlisting}[style=codeblock,caption={\code{src/UseCases/gfmc/run.jl}}]
function _reconfigure!(sim::GFMCSim)
    indices, mean_weight = resample_indices(sim.reconfiguration, sim.rng, sim.weights)
    copy_columns!(sim.backend, sim.proposal_positions, sim.positions, indices)
    copy_batch_data!(sim.backend, sim.proposal_data, sim.state_data, indices)

    tmp_positions = sim.positions
    sim.positions = sim.proposal_positions
    sim.proposal_positions = tmp_positions

    tmp_data = sim.state_data
    sim.state_data = sim.proposal_data
    sim.proposal_data = tmp_data

    fill_ones!(sim.backend, sim.weights)
    return mean_weight
end
\end{lstlisting}

So instead of explicit birth/death at every step, this code accumulates weights
for several steps and only then converts weight imbalance into ancestry counts.
That is why this implementation is best described as a weighted fixed-population
GFMC projector.

\section{Reference-Energy Control}

The reference energy update is also fixed-population in spirit:

\begin{lstlisting}[style=codeblock,caption={\code{src/UseCases/gfmc/run.jl}}]
function _update_reference_energy!(sim::GFMCSim, mean_weight::Real)
    window = sim.params.energy_window
    if isempty(sim.energy_mean_history)
        Eblock = estimate_energy(sim)
    else
        hi = lastindex(sim.energy_mean_history)
        lo = max(1, hi - window + 1)
        block = @view sim.energy_mean_history[lo:hi]
        Eblock = sum(block) / length(block)
    end

    dt_block = sim.params.dt * sim.params.reconfiguration_interval
    correction = sim.params.feedback == 0 ? 0.0 :
        (sim.params.feedback / dt_block) * log(max(Float64(mean_weight), eps(Float64)))
    sim.ET = Eblock - correction
    return sim
end
\end{lstlisting}

In equations, the update is
\begin{equation}
\ET \leftarrow E_{\rm block} -
\frac{\texttt{feedback}}{\dtau \cdot \texttt{reconfiguration\_interval}}
\log(\overline w),
\end{equation}
where $\overline w$ is the mean weight returned by the reconfiguration stage.

Interpretation:
\begin{itemize}
    \item if the ensemble weights are trending upward, the correction lowers
    future growth,
    \item if the ensemble weights are trending downward, the correction raises
    future growth.
\end{itemize}

This plays the role that population control plays in ordinary DMC, but it is
adapted to a constant-walker ensemble.

\section{Energy Estimation}

The measured energy is the weighted mean of the current branch-energy field:

\begin{lstlisting}[style=codeblock,caption={\code{src/UseCases/gfmc/sim.jl}}]
function estimate_energy(sim::GFMCSim)::Float64
    numerator = 0.0
    denom = 0.0
    @inbounds for j in eachindex(sim.weights)
        w = sim.weights[j]
        numerator += w * sim.state_data.branch_energy[j]
        denom += w
    end
    return denom > 0 ? numerator / denom : 0.0
end
\end{lstlisting}

So the estimator is
\begin{equation}
E \approx \frac{\sum_j w_j \EL(R_j)}{\sum_j w_j},
\end{equation}
which is exactly the mixed-estimator logic you expect from importance-sampled
projector Monte Carlo.

\section{How You Use the Engine}

\subsection{Generic Hamiltonian path}

The generic path is:
\begin{enumerate}
    \item define a \code{Hamiltonian},
    \item define a \code{TrialWF} if you want importance sampling,
    \item wrap it as \code{ImportanceGuiding(trial, H)},
    \item create \code{GFMCParams},
    \item create \code{GFMCSim},
    \item call \code{run\_gfmc!}.
\end{enumerate}

\begin{lstlisting}[style=codeblock,caption={Minimal generic usage from \code{docs/GFMC_USER_GUIDE.md}}]
using Random
using System1D

bc = PeriodicBoundary1D(0.0, 1.0)
H = Hamiltonian(1, 0.5, R -> 0.5 * sum(abs2, R), bc)
trial = TrialWF(
    R -> -0.5 * sum(abs2, R),
    R -> -R,
    R -> -length(R),
)
guiding = ImportanceGuiding(trial, H)

params = GFMCParams(0.02, 100, 20, 256, 0.5, 0.2, 5, 5.0, 20)
walkers = [Walker([rand()]) for _ in 1:params.targetN]

sim = GFMCSim(H, params, walkers, MersenneTwister(7);
    guiding=guiding,
    nodepolicy=NoNode(),
    reconfiguration=SystematicReconfiguration(),
)

run_gfmc!(sim; snapshot_steps=[0, params.nsteps], debug=true)
\end{lstlisting}

\subsection{VMC warm start}

The repository also includes a one-call helper that first samples
$|\Psi_T|^2$ with VMC and then launches GFMC:

\begin{lstlisting}[style=codeblock,caption={Excerpt from \code{src/UseCases/common/warm_start.jl}}]
function run_gfmc_with_vmc_init(
    h,
    gfmc_params::GFMCParams,
    init_configs,
    trial_wf,
    vmc_params::VMCParams;
    vmc_rng::AbstractRNG=Random.default_rng(),
    gfmc_rng::AbstractRNG=Random.default_rng(),
    proposal::Union{AbstractVMCProposal,Symbol}=DriftGaussianProposal(),
    guiding::AbstractGuiding=ImportanceGuiding(trial_wf, h),
    nodepolicy::AbstractNodePolicy=NoNode(),
    backend::AbstractGFMCBackend=CPUBackend(),
    reconfiguration::AbstractGFMCReconfiguration=SystematicReconfiguration(),
    ...
)
    positions = sample_vmc_positions(...)
    sim = GFMCSim(h, gfmc_params, positions, gfmc_rng; ...)
    return run_gfmc!(sim; ...)
end
\end{lstlisting}

This helper is conceptually clean:
\begin{itemize}
    \item VMC provides an initial ensemble closer to $|\Psi_T|^2$,
    \item GFMC then projects that ensemble in imaginary time.
\end{itemize}

\subsection{Specialized fermion-ring path}

For the ring model, you can bypass the generic trial/guiding setup because the
kernel knows the model's built-in trial structure:

\begin{lstlisting}[style=codeblock,caption={Minimal ring-model usage from \code{docs/GFMC_USER_GUIDE.md}}]
using Random
using System1D

model = SpinlessFermionRing1D(3, 1.0, 3.0, 1.0;
    lambda=-0.5,
    node_tol=1e-8,
    trig_eps=1e-10,
)
params = GFMCParams(0.003, 50, 10, 128, -0.2, 0.1, 2, 5.0, 10)
positions = sample_uniform_configurations(model, params.targetN, MersenneTwister(101))

sim = GFMCSim(model, params, positions, MersenneTwister(202);
    use_guiding=true,
    nodepolicy=FixedNode(),
)

run_gfmc!(sim; snapshot_steps=[0, params.nsteps])
\end{lstlisting}

\section{What Is Most Important to Remember}

If you want one precise sentence that matches the code, it is this:

\begin{quote}
This GFMC implementation approximates the repeated action of the
imaginary-time propagator by alternating importance-sampled drift-diffusion
proposals, Metropolis correction, multiplicative short-time weights, and
periodic fixed-population reconfiguration.
\end{quote}

The most important caveats are:
\begin{itemize}
    \item it is a fixed-population weighted projector, not explicit branching
    DMC,
    \item unguided paths set \code{sign = 1.0}, so fixed-node logic only has
    real effect when the active kernel supplies sign information,
    \item the recorded energy is the weighted branch-energy estimator, not a
    descendant-weight or forward-walking estimator,
    \item the GFMC path stores real coordinates only and does not propagate the
    phase factors associated with twisted-boundary wrapping.
\end{itemize}

\section{How to Compile This File}

From the repository root, run:

\begin{lstlisting}[style=codeblock]
mkdir -p docs/build
latexmk -pdf -output-directory=docs/build docs/GFMC_CODE_WALKTHROUGH.tex
\end{lstlisting}

The expected PDF output path is:

\begin{lstlisting}[style=codeblock]
docs/build/GFMC_CODE_WALKTHROUGH.pdf
\end{lstlisting}

If you prefer plain \code{pdflatex}, this also works:

\begin{lstlisting}[style=codeblock]
mkdir -p docs/build
pdflatex -interaction=nonstopmode -halt-on-error \
  -output-directory=docs/build docs/GFMC_CODE_WALKTHROUGH.tex
\end{lstlisting}

\end{document}
